% Pasted excerpt, user-supplied 2026-08-25 (source paper/section not otherwise
% identified in this repo -- referred to by the user as "Owen_section_D";
% likely a draft/section from a follow-on or companion paper to
% brower_2018_lattice_phi4_riemann_S2.pdf, not yet vendored as a full PDF
% here). Kept verbatim (including \rev{}/\begingroup\color{violet} revision
% markup) rather than cleaned up, so it can be diffed against the eventual
% published version if one is added later.
%
% Content: "Agreement with the Ising CFT" -- defines the Legendre-projected
% lattice two-point-function coefficients F_l(a), their continuum closed
% form F_l^cont(Delta) (same recursion as Brower 2018 Eq 5.2-5.3, different
% overall normalization convention), a dimension-independent conformal
% symmetry-breaking diagnostic delta_l(a) built purely from F_0, F_1,
% F_{l-1}, F_l (no need to know Delta_sigma), and a direct scaling-dimension
% estimator Delta_s(a) = 2*F_1/(F_1+F_0). See IsingS2_precision journal.md
% 2026-08-25 for how this compares to Brower Sec. 5.1 and the campaign's
% existing M_l[m,m'] cumulant-based symmetry test.

\subsection{Agreement with the Ising CFT}
\label{sec:CFTtest}

Now that we have confirmed that the spin-spin correlation function in our lattice theory becomes rotationally symmetric in the continuum limit, we would like to check that it also agrees with the exactly known analytical result for the 2d Ising CFT on a 2-sphere:
\begin{equation}
\label{eq:CFT_sphere_2pt}
\rev{\< \sigma(\hat n_1)  \sigma(\hat n_2)\>=\dfrac{1}{|\hat n_1 -\hat n_2|^{2\Delta_{\sigma}}}} = \sum_{\ell} \dfrac{2 \ell + 1}{2} F^{cont}_{\ell} P_{\ell} ( \rev{z})
\end{equation}
is a Legendre expansion of the correlator
for scalar primary   in the relative angle, $ z= \cos( \theta)$. The normalization is chosen to match the UV
divergence, $1/|x_1 - x_2|^{2 \Delta}$, on the flat tangent \rev{plane}.

\begingroup\color{violet}
This is not an independent input, but follows from conformal covariance
under the stereographic projection relating $\mS^2$ to $\mR^2$: a primary
operator transforms as
$\sigma_{\mS^2}(\hat n) = \Omega(\hat n)^{-\Delta_\sigma}\, \sigma_{\mR^2}(x)$,
where $\Omega$ is the conformal factor of the map~\cite{Rychkov:2016iqz}.
The flat two-point function
$\langle \sigma(x_1) \sigma(x_2) \rangle = 1/|x_1 - x_2|^{2\Delta_\sigma}$ is
therefore mapped exactly onto the chordal-distance form in
Eq.~(\ref{eq:CFT_sphere_2pt}) for a sphere of any radius $R$, with $R$
entering only as an overall normalization that cancels in the
scale-invariant ratio $\Delta_\sigma = 2 F_1/(F_1 + F_0)$ measured below.
At fixed physical radius $R$, the finite radius is part of the target CFT
geometry rather than an unwanted periodic-image effect, so no separate
infinite-volume extrapolation is required.  The relevant hierarchy is
$a\ll R$, and the continuum limit is taken as $a/R\rightarrow0$ (equivalently
$N\rightarrow\infty$ at fixed $R$).  Although a finite lattice on the compact
sphere admits no sharp thermodynamic phase transition, the geometric
couplings define a sequence that approaches the renormalization-group fixed
point established in Sec.~\ref{sec:WMcontinuum}; they are not claimed to give
an exactly massless system at finite refinement.  The residual cutoff and
mass effects are $O(a^2)$, and Eq.~(\ref{eq:CFT_sphere_2pt}) becomes exact only
in this fine-lattice limit.
\endgroup

 The  coefficients
\begin{equation}
 F^{cont}_l = \int^1_{-1}dz \dfrac{1}{(2 - 2 z)^\Delta} P_l(z)
 \end{equation}
obey the \rev{recursion} relation
\begin{equation}
\label{eq:2pt_coeff}
    F^{cont}_{l} = \dfrac{l - 1 +\Delta }{l + 1 - \Delta  }   F^{cont}_{l-1}\quad,\quad F^{cont}_0 = \rev{\dfrac{2^{1-2\Delta}}{1-\Delta}}
\end{equation}
\begingroup\color{violet}
The corresponding lattice coefficients $F_\ell(a)$ are measured
directly from the configurations generated by our Monte Carlo
simulations, as described in Appendix~\ref{app:MCmeasure}.
To check that the lattice simulation agrees with the analytic result,
we must compute these coefficients at finite lattice spacing and then
take the continuum limit. The expression is
\endgroup
\bea
F_l(a) &=&\dfrac{1}{2 \pi  }\sum_{ij}  \sqrt{g_i g_j}\left(\sum^{m=l}_{m=-l} \rev{Y^*_{l,m}(\hat r_i)Y_{l,m}(\hat r_j)}\right)\< s_i s_j \> \nn
&=&\dfrac{2l + 1}{2 \pi \times 4\pi }\sum_{ij}  \sqrt{g_i g_j} P_l(\hat r_i \cdot \rev{\hat r_j})  \< s_i s_j \>
  \eea
  using the addition theorem and integration rule  in Appendix A, where
  \rev{it is shown that this reproduces the} CFT correlator $\< \sigma_1 \sigma_2 \>$\rev{ in} the continuum limit.
  \begingroup\color{violet}
  The Ising correlator, $\< s_i s_j \>$, is computed by Monte Carlo
  simulation at each of the lattice refinements shown in
  Figs.~\ref{fig:cft_break_mod} and \ref{fig:sphere_delta_sigma}.
  This is then extrapolated  to the continuum ($a/r \rightarrow  0$) and compared with the exact CFT result up to a
  multiplicative  renormalization factor relating the Ising lattice
  spin $s_i$  to the primary operator   $\sigma(r)$. This does not affect
  homogeneous ratios of the $F_\ell$, but it does fix the overall
  normalization through
  \endgroup
\be
F_0(a)=\dfrac{ 1}{2 \pi \times 4\pi }\sum_{ij}  \sqrt{g_i g_j}   \< s_i s_j \> \rightarrow F_0(0) = 2\< M^2 \>
\ee
where $\sum_i\sqrt{g_i} s_i/4 \pi \rightarrow M$ is the total \rev{magnetization} operator
  in the continuum \rev{extrapolation.}
  \begingroup\color{violet}
  Equivalently, defining
  $M(a)=\sum_i\sqrt{g_i}s_i/(4\pi)$ at finite lattice spacing gives
  $F_0(a)=2\langle M(a)^2\rangle$ directly, so the normalization is
  fixed by the measured magnetization variance rather than by a fit.
  \endgroup

  Now \rev{rewriting} the recursion  relation  as
  \be
  l(F^{cont}_{l}  - F^{cont}_{l-1}) = (\Delta -1)\rev{(}F^{cont}_{l} + F^{cont}_{l-1} )
\ee
\rev{and}
 \rev{eliminating} $\Delta_{s}$ we find a
test for a conformal symmetry breaking measurement \rev{independent}
of the \rev{dimension}.
\begin{equation}
    \delta \ell(a) = 1 - \dfrac{(F_0 - F_1)(F_{\ell-1} + F_{\ell})}{\ell (F_0 + F_1)(F_{\ell-1} - F_{\ell})}
\end{equation}
which should go to zero in the continuum limit. This quantity is plotted as a function of lattice spacing in Fig.~\ref{fig:cft_break_mod}, showing the expected behavior for both the octahedral and icosahedral lattice for all \rev{irreps} of O(3) up to $\ell = 12$.
\begin{figure*}[ht]
   % \centering
    \includegraphics[width=0.49\textwidth]{q5_cft_break.png}
    \includegraphics[width=0.49\textwidth]{q4_cft_break.png}
    \caption{Conformal symmetry breaking measurement for the spin-spin two-point function as a function of lattice spacing using the modified discretization of the sphere comparing our  icosahedral (left) lattice to the octahedral (right)  lattice up to a refinement of $L = 64$.}
    \label{fig:cft_break_mod}
\end{figure*}
We can also measure the scaling exponent of the $s$ operator on the lattice via
\begin{equation}
    \Delta_{s}(a) = \dfrac{2 F_1}{F_1 + F_0}
\end{equation}
which is plotted as a function of lattice spacing in Fig. \ref{fig:sphere_delta_sigma}. Extrapolating to the continuum limit by fitting to a quadratic polynomial in $a$, we obtain the results $\Delta_{s} = 0.125048(44)$ with $\chi^2/\text{dof}=1.8$ for the octahedral lattice and $\Delta_{s} = 0.124985(47)$ with $\chi^2/\text{dof}=1.8$ for the icosahedral lattice, both in excellent agreement with the exact value of $1/8$ and with a relative uncertainty of about $0.03\%$.
\begin{figure*}[ht]
%    \centering
    \includegraphics[width=0.49\textwidth]{q5_delta_sigma.png}
    \includegraphics[width=0.49\textwidth]{q4_delta_sigma.png}
    \caption{Continuum extrapolation of the scaling exponent $\Delta_{s}$ for the modified icosahedral (left) lattice compared to modified octahedral (right). }
    \label{fig:sphere_delta_sigma}
\end{figure*}

Our results indicate that by using the modified discretization of $S^2$, our lattice simulation accurately captures the properties of the lowest $\mathbb{Z}_2$-odd operator $s$. The 2d Ising CFT of course contains an infinite set of operators related by the Virasoro algebra. Measurement of additional operators using our lattice model is a subject for future study outside the present scope, however based on the high level of accuracy of our results and the delicate nature of a properly-tuned conformal field theory, we believe it is likely that our lattice formulation is sufficient to capture the properties of all of the operators in the 2d Ising CFT in the continuum limit, given sufficient computational resources.


Another approach taken recently is to use the so-called ``fuzzy sphere'' which projects the states of the theory in the basis of spherical harmonics instead of using a traditional spatial lattice~\cite{Zhu2022UncoveringCS,Hu2023OperatorPE}. In this case, the lattice cutoff is replaced by a maximum spherical harmonic quantum number, which is a more manifest way to ensure that rotational symmetry is restored in the continuum limit. This method shows promising results for the lowest operators of the 3d Ising CFT, though it is unclear how difficult it would be to generalize to other theories.
